\subsection{THETA3-C33: the higher-CM graph is a labelled horizontal interior special-cycle closure}\label{C33}

The remaining formal comparison in C31 can be made on its
specified good-prime toroidal chart. The result concerns the
labelled horizontal closure. It does not identify additional
vertical components of an integral special-homomorphism
moduli space, nor a correspondence from the original compact
ball to this larger PEL model.

\paragraph{Statement and the precise closure.}
Keep all the data and prime-to-$p$ choices of C31. Write
$A=B^3$ for the abelian part at the cusp, and let
$G^{\natural}$ be its Raynaud extension by the cusp torus.
The associated degeneration datum also contains a lattice
$Y$, a period map $\iota:Y_\eta\to G^{\natural}_\eta$,
and the polarization and level structures. Work on the
finite \'etale formal pullback of a smooth toroidal chart
which labels the kernel of $[m]\varphi$ in C31. Finite
\'etale covers of its closed stratum extend uniquely to
its formal completion, so this is an actual chart of the
same formal PEL model.
This is the finite \'etale equivalence for henselian
pairs, Stacks Project, Lemma~15.13.2; at each
nilpotent stage it is Lemma~15.11.2, Tag~0ALI.

For one of C31's columns
\[
 h_+=(\bar a,1,-1),\qquad
 h_1=(\bar a,0,-1),\qquad h_2=(0,1,-1):B\longrightarrow A,
                                                        \tag{C33a}
\]
select the zero-labelled locus where $h^\vee c=0$.
It is precisely the corresponding graph $\Gamma_j$ in
C31, pulled back in the toroidal directions. Then this
formal graph is the scheme-theoretic closure of the
labelled characteristic-zero interior special-homomorphism
locus. Its Hermitian norm is, respectively,
\[
              a\bar a+2,\qquad a\bar a+1,\qquad 2.
                                                        \tag{C33b}
\]
In particular C31's graph contraction expresses the actual
Poincar\'e class using these labelled horizontal interior
closures, not just abstract boundary loci.

To specify ``horizontal'' without ambiguity, choose local
smooth cone coordinates $q_1,\ldots,q_r$, and put
$q=q_1\cdots q_r$. The abelian interior is obtained by
inverting $q$. In a mixed-characteristic chart we take
the closure from the locus where both $p$ and $q$ are
invertible. Thus the closure ideal is saturated by $pq$.
If a chosen integral interior component is already its
horizontal closure, one may equivalently take its
$q$-saturated closure. We do not assert that arbitrary
$p$-vertical components disappear under $q$-saturation.

\paragraph{The exact fully faithful theorem.}
The source is Lan's corrected author revision of
\emph{Arithmetic compactifications of PEL-type Shimura
varieties}, Theorem~4.4.16, printed pp.~254--255, and
especially Corollary~4.5.5.5 with its complete proof,
p.~303. The latter explicitly treats an arbitrary
homomorphism of generic abelian varieties: it gives
uniquely a morphism of degeneration data
\[
 (f_Y,f^{\natural}):(Y_1,G_1^{\natural})
                     \longrightarrow(Y_2,G_2^{\natural}),
 \qquad f^{\natural}\iota_1=\iota_2f_Y.                 \tag{C33c}
\]
It is not restricted to polarization-preserving
isomorphisms. Its hypotheses, stated in Section~4.1,
printed pp.~207--208, are a noetherian normal domain
complete for an ideal, with the indicated semiabelian
and isotrivial-torus hypotheses. The normal completed
charts below satisfy them. Corollary~4.5.5.5 identifies
its assertion with Faltings--Chai, Chapter~III,
Theorem~7.1. We use the corrected author's statement
and proof, not an inference from an equivalence of
polarized groupoids alone.

The compatibility with duality and polarization is
explicit in Lan's Construction~4.5.4.27 and
Lemma~4.5.4.28, printed pp.~299--300. The actual PEL
formal family is the Mumford family under Lan's
\emph{Toroidal compactifications of PEL-type Kuga
families}, Theorem~1.41(5), pp.~16--17. Hence applying
(C33c) compares to the actual interior abelian family
on that chart.

\paragraph{There is no extra source-lattice condition.}
Let $S=\operatorname{Spec}(R')$ be one of the normal
complete chart bases in the preceding theorem. The
good-reduction source $B/S$ has degeneration datum
\[
                         [0\longrightarrow B].
\]
A morphism from it to $[Y\longrightarrow G^{\natural}]$
therefore consists exactly of a group homomorphism
$\widetilde h:B\to G^{\natural}$: the lattice map
$0\to Y$ is zero, and the compatibility in (C33c)
is vacuous. Such a lift of $h:B\to A$ exists if
and only if the pulled-back torus extension splits.
By C31 this is precisely
\[
                              h^\vee c=0.              \tag{C33d}
\]
The splitting is unique because $\Hom(B,T)=0$.
It is automatically $\mathcal O_E$-linear by uniqueness.
Thus the fully faithful theorem yields a unique actual
homomorphism $f:B\to G$ of degenerating group schemes,
and in particular its interior abelian-fibre
homomorphism, exactly on (C33d).

For clarity the dual period condition gives nothing
additional either. The dual datum has the form
$[X\to G^{\vee,\natural}]$, with abelian quotient
$A^\vee$ and projected period map $c:X\to A^\vee$.
The dual map to $[0\to B^\vee]$ has semiabelian
extension part $h^\vee\circ\pi^\vee$ and zero lattice
map. Its period condition is exactly $h^\vee c=0$,
again (C33d). Consequently the target's period
trivialization and toroidal $q$ parameters do not
supply a further equation for this source. The
polarization relation between $c$ and $c^\vee$ is
retained in the given target degeneration datum;
it is not discarded or separately postulated.

\paragraph{The norm is exactly the abelian-part norm.}
Functoriality, duality and the polarization construction
identify the abelian part of $f^\vee\lambda_G f$ with
$h^\vee\lambda_Ah$. Its source and target both have
zero torus and lattice parts, so this determines the
whole homomorphism of good abelian schemes. Thus
\[
       \lambda_B^{-1}f^\vee\lambda_Gf
          =\lambda_B^{-1}h^\vee\lambda_Ah.              \tag{C33e}
\]
Here the inverse polarization is an equality of
quasi-homomorphisms; its right side in the three
cases is the actual integral endomorphism in
(C33b). Indeed $\lambda_A$ is the product
polarization and its Rosati involution sends
$a$ to $\bar a$, so the three matrix products
are $a\bar a+1+1$, $a\bar a+1$, and $1+1$.
All are totally positive. No factor from the
period lattice, the torus, or the prime-to-$p$
boundary isogeny changes these norms. The
isogeny changes the extension parameter as
recorded in C31, and its zero-labelled graph
has already accounted for that change.

\paragraph{Scheme structure of the horizontal closure.}
We give the argument for one graph; it applies to
all three. Let $\mathfrak S$ be the selected formal
chart. \'Etale locally its toroidal direction is a
smooth formal torus embedding over the boundary
abelian parameter scheme. The graph $\mathfrak Z$
is the inverse image of $\Gamma_j$. Its equation
is independent of those directions. Therefore
$\mathfrak Z$ is smooth over the original
mixed-characteristic base, is normal and flat,
and multiplication by each $q_i$ is injective
on its local rings. In particular its ideal
$I_{\mathfrak Z}$ is saturated for both $p$
and $q$.

On each normal complete local ring of
$\mathfrak Z$, (C33d) splits the extension,
so (C33c) constructs the actual interior
homomorphism with norm (C33b). Uniqueness
makes these homomorphisms agree on overlaps.
This gives the inclusion of $\mathfrak Z$
in the closure of the selected interior
locus: its open part with $pq$ invertible
is schematically dense.

We verify the converse with the actual
scheme structures. In characteristic zero
the labelled special-homomorphism moduli
locus is smooth after neat level. Here is
an infinitesimal justification, which is
stronger than a description of its set
of complex points. Write $t=f^\dagger f$
for the totally positive element (C33b),
where $f^\dagger=\lambda_B^{-1}f^\vee\lambda_G$.
It is invertible in $E$, and
\[
                  e=f t^{-1}f^\dagger
\]
is a self-adjoint idempotent in
$\operatorname{End}^0_{\mathcal O_E}(G)$.
It splits the target up to isogeny into
the fixed factor $B$ and its orthogonal
abelian complement $C$. This assertion
works over arbitrary local Artinian
characteristic-zero test rings: clear
one integer denominator, and take the
images of the two complementary integral
endomorphisms. Equivalently take their
connected kernels. The resulting product
map is an isogeny, and its finite kernel
is \'etale because every nonzero integer
is invertible on the test ring.
The image and connected-kernel construction
is valid over these Artinian bases: the
formal logarithm in characteristic zero
identifies the completed group law with
its Lie algebra, where the idempotent is
an actual projection. Its kernel is
formally smooth of the constant required
rank; the connected proper kernel is
therefore an abelian scheme.

After one finite \'etale torsion level,
that kernel and all its action, isotropy
and level constraints have unique lifts
across every nilpotent thickening. The
polarization on the first factor is the
fixed polarization $\lambda_B t$; its
mixed term with $C$ is zero because
$e$ is self-adjoint. Consequently the
deformation functor of the target
together with the labelled map $f$
is exactly the deformation functor of
the complementary PEL abelian variety,
with those discrete isogeny data fixed.
Conversely, deforming the complement
and taking the quotient by the uniquely
lifted finite kernel reconstructs the
target, its polarization and $f$.
These two constructions are inverse
on Artinian test rings and on their
isomorphisms.

The complementary PEL functor is smooth
in characteristic zero. In the orientation
selected by the source CM type, its
signatures are $(3,1)$ instead of
$(4,1)$ at each real place; the conjugate
orientation reverses both pairs. This
is the ordinary smooth unitary PEL
moduli problem, also supplied by Lan's
Theorem~1.8 in the cited Kuga-family
paper. Thus the actual labelled locus
is smooth and reduced, of codimension
$d=[E^+:\Q]$. The unramifiedness of the
Hom functor by rigidity is used only
to choose its local labelled closed
immersion; it is not used to deduce
reducedness. This last localization is
the \'etale local structure of unramified
morphisms, Stacks Project, Lemma~41.17.1,
Tag~04HH.

Its scheme-theoretic horizontal closure
is reduced: locally its ideal is the
contraction of a radical ideal after
inverting $pq$, which is radical.
The rings in question are excellent;
completion along the boundary preserves
reducedness, and normalization is finite.
For exact algebra references, completion
of a G-ring is a regular map by Stacks
Project, Lemma~15.51.14, Tag~0AH2;
reducedness and normality ascend under
such a map by Lemmas~10.163.6 and
10.163.8, Tags~0C21 and~0C22.
Finite normalization follows from
Lemma~10.162.10(4), Tag~032Y, or from
the Nagata property of excellent rings.
Normalize a component of this closure
which meets the selected formal chart,
and complete at a boundary point. After
choosing a local domain component this
is a normal complete base to which
(C33c) applies. Its generic interior
homomorphism extends uniquely, and its
Raynaud abelian-part homomorphism is the
specified label $h$. Indeed equality
with that boundary label holds on the
closed fibre, and rigidity gives equality
on every Artinian thickening; separatedness
and completeness give equality over the
complete base. This is the same
normalized homomorphism-rigidity argument
used in C23. Consequently (C33d)
holds identically on that normal component.
It therefore maps into $\mathfrak Z$.
All components of the reduced closure
satisfy this, so the closure ideal contains
$I_{\mathfrak Z}$. Together with the first
inclusion this proves equality of the
reduced closed formal subschemes. Equivalently,
their ideals agree after inverting $pq$,
and both are $pq$-saturated.

This proof retains the label: other finite
kernel labels give their corresponding
translated graphs, and other special
homomorphisms are not identified with the
zero label. It also retains the horizontal
qualification. In particular it gives no
formula for a derived vertical intersection
length, and does not change the sign or
coherence constraints in C06 and the
arithmetic Siegel--Weil comparisons.

\paragraph{Consequences for the integral kernel formula.}
Let $i$ denote the inclusion of the
closed boundary stratum. The $q_i$
form a regular sequence on every
closure just constructed, so its
refined boundary pullback is exactly
$[\Gamma_j]$, with no excess term.
Substitute these refined restrictions
for $\Gamma_{j_+},\Gamma_{j_1},\Gamma_{j_2}$
in (C31c). The actual integral Chow
identity and polarization contraction
therefore remain an identity on the
boundary abelian parameter scheme,
now expressed as restrictions of the
specified interior horizontal closures.
This does not assume a proper projection
forgetting $T$ on the entire toroidal
chart. The C23/A6 finite CM
projector therefore uses specified
labelled interior horizontal special
cycles in this larger unitary model,
with the exact prime-to-$p$ descent
indices already retained. This closes
the degeneration-data condition left
open in C31. An interior correspondence
from the original compact 032 ball,
and a comparison of vertical derived
intersection numbers, remain additional
geometric requirements.

\framework{The lifting obstruction is the
Poincar\'e extension class from SGA~7~I,
VII~2.9.5 and VIII~3.1--3.2, as in C31.
The new input here is the fully faithful
Mumford degeneration theorem for arbitrary
homomorphisms, Lan Corollary~4.5.5.5,
p.~303. BHKRY Remark~3.6.2, printed
p.~53, describes the elliptic instance
as homomorphisms into the universal
one-motive with its Hermitian norm;
Theorem~3.7.1(8), p.~54, identifies its
formal closures. The proof above keeps
the higher-CM source and proves the
labelled horizontal comparison in the
specified chart.}
